%
% Chapter 10: Training Process and Frameworks

\clearpage
\cleardoublepage

\chapter{Training Process and Frameworks}

\section{The Training Loop}

Neural network training consists of iterative optimization over data.

\begin{algorithm}[Training Loop]
\begin{enumerate}
    \item Initialize model weights
    \item For each epoch:
    \begin{enumerate}
        \item Shuffle training data
        \item For each batch:
        \begin{enumerate}
            \item Forward pass: compute predictions
            \item Compute loss
            \item Backward pass: compute gradients
            \item Update weights using optimizer
        \end{enumerate}
        \item Evaluate on validation set
        \item Save checkpoint if improved
    \end{enumerate}
    \item Return best model
\end{enumerate}
\end{algorithm}

\section{Hyperparameters}

Critical hyperparameters that affect training:

\subsection{Learning Rate}

The most important hyperparameter.

\begin{itemize}
    \item \textbf{Too high:} Divergence, exploding loss
    \item \textbf{Too low:} Slow convergence, local minima
    \item \textbf{Typical range:} $10^{-5}$ to $0.1$
    \item \textbf{Default for Adam:} $10^{-3}$
    \item \textbf{Default for SGD:} $10^{-2}$
\end{itemize}

Finding good learning rates:
\begin{itemize}
    \item Learning rate range test (LR finder)
    \item Start with known good defaults
    \item Reduce by factor of 10 if loss diverges
\end{itemize}

\subsection{Batch Size}

Number of samples per training step.

\begin{itemize}
    \item \textbf{Small batch:} Noisier gradients, better generalization
    \item \textbf{Large batch:} Smoother gradients, faster but needs LR scaling
    \item \textbf{Typical values:} 16, 32, 64, 128, 256
    \item \textbf{Large batch scaling:} $\alpha_{large} = \alpha_{small} \times \frac{B_{large}}{B_{small}}$
\end{itemize}

\begin{note}[Batch Size vs. Generalization]
Research shows smaller batches often achieve better generalization, but the effect depends on other factors like learning rate and network architecture.
\end{note}

\begin{table}[h]
    \centering
    \caption{Hyperparameter Ranges}
    \begin{tabular}{L{2.5cm} L{3cm} L{3.5cm}}
        \toprule
        \textbf{Hyperparameter} & \textbf{Typical Range} & \textbf{Notes} \\
        \midrule
        Learning Rate & $10^{-5}$ to $0.1$ & Adam: $10^{-3}$, SGD: $10^{-2}$ \\
        Batch Size & 16-512 & Depends on GPU memory \\
        Epochs & 10-1000 & Depends on dataset size \\
        Weight Decay & $10^{-6}$ to $10^{-2}$ & AdamW: $10^{-2}$ typical \\
        Dropout & 0.1-0.5 & Depends on regularization need \\
        Momentum & 0.9-0.999 & For SGD \\
        \bottomrule
    \end{tabular}
\end{table}

\section{Regularization}

Techniques to prevent overfitting:

\subsection{Dropout}

Randomly zero out activations during training:
\begin{equation}
\mathbf{y} = \mathbf{W} \cdot (\mathbf{x} \odot \mathbf{m})
\end{equation}
where $\mathbf{m} \sim \text{Bernoulli}(p)$ is the mask.

During inference, no dropout (or use inverted dropout).

\begin{properties}[Dropout Effects]
\begin{itemize}
    \item Ensemble of sub-networks
    \item Reduces co-adaptation
    \item Increases training time slightly
    \item Typical dropout rate: 0.1-0.5
\end{itemize}
\end{properties}

\subsection{Weight Decay}

L2 regularization:
\begin{equation}
\mathcal{L}_{reg} = \mathcal{L} + \frac{\lambda}{2}\|\mathbf{W}\|_F^2
\end{equation}

Equivalent to:
\begin{equation}
\mathbf{W} \leftarrow (1 - \alpha\lambda)\mathbf{W} - \alpha \nabla \mathcal{L}
\end{equation}

For AdamW: Weight decay is applied independently of the adaptive learning rate.

\section{Early Stopping}

Stop training when validation performance stops improving:
\begin{enumerate}
    \item Monitor validation loss
    \item Save best model
    \item Stop if no improvement for $N$ epochs
    \item Return best model
\end{enumerate}

Typical patience: 10-20 epochs

\section{Model Checkpointing}

Save model state during training:
\begin{itemize}
    \item \textbf{Architecture:} Model definition
    \item \textbf{Weights:} Learned parameters
    \item \textbf{Optimizer state:} Momentum, etc.
    \item \textbf{Epoch:} Current training progress
    \item \textbf{Best metric:} For model selection
\end{itemize}

\section{Backpropagation}

The algorithm for computing gradients through the computational graph.

\begin{algorithm}[Backpropagation]
\begin{enumerate}
    \item Forward pass: compute all intermediate values
    \item Compute loss $\mathcal{L}$
    \item Backward pass:
    \begin{enumerate}
        \item $\frac{\partial \mathcal{L}}{\partial \mathbf{y}}$ at output
        \item For each layer (reverse order):
        \begin{enumerate}
            \item Compute $\frac{\partial \mathcal{L}}{\partial \mathbf{x}} = \frac{\partial \mathcal{L}}{\partial \mathbf{y}} \cdot \frac{\partial \mathbf{y}}{\partial \mathbf{x}}$
            \item Compute $\frac{\partial \mathcal{L}}{\partial \mathbf{W}}$ for weight update
        \end{enumerate}
    \end{enumerate}
\end{enumerate}
\end{algorithm}

\begin{note}[Chain Rule]
Backpropagation applies the chain rule repeatedly:
\begin{equation}
\frac{\partial \mathcal{L}}{\partial x} = \frac{\partial \mathcal{L}}{\partial y} \cdot \frac{\partial y}{\partial x}
\end{equation}
For complex graphs, this is computed efficiently using automatic differentiation.
\end{note}

\section{PyTorch Framework}

PyTorch is the most popular framework for deep learning research.

\begin{code}
    \captionof{listing}{Training Loop in PyTorch}
    \label{code:training_pytorch}
\begin{minted}{python}
import torch
import torch.nn as nn
import torch.optim as optim
from torch.utils.data import DataLoader

# Model
model = nn.Sequential(
    nn.Linear(784, 256),
    nn.ReLU(),
    nn.Dropout(0.2),
    nn.Linear(256, 10)
)

# Loss and optimizer
criterion = nn.CrossEntropyLoss()
optimizer = optim.Adam(model.parameters(), lr=0.001, weight_decay=0.01)

# Training loop
device = torch.device('cuda' if torch.cuda.is_available() else 'cpu')
model = model.to(device)

for epoch in range(num_epochs):
    model.train()
    train_loss = 0.0
    for batch_idx, (data, target) in enumerate(train_loader):
        data, target = data.to(device), target.to(device)
        
        optimizer.zero_grad()
        output = model(data.view(data.size(0), -1))
        loss = criterion(output, target)
        loss.backward()
        optimizer.step()
        
        train_loss += loss.item()
    
    # Validation
    model.eval()
    val_loss = 0.0
    correct = 0
    with torch.no_grad():
        for data, target in val_loader:
            data, target = data.to(device), target.to(device)
            output = model(data.view(data.size(0), -1))
            val_loss += criterion(output, target).item()
            pred = output.argmax(dim=1)
            correct += pred.eq(target).sum().item()
    
    print(f'Epoch {epoch}: Train Loss={train_loss/len(train_loader):.4f}, '
          f'Val Loss={val_loss/len(val_loader):.4f}, '
          f'Accuracy={100.*correct/len(val_loader.dataset):.2f}%')
\end{minted}
\end{code}

\section{CUDA and GPU Training}

PyTorch uses CUDA for GPU acceleration.

\begin{code}
    \captionof{listing}{CUDA Operations in PyTorch}
    \label{code:cuda_pytorch}
\begin{minted}{python}
import torch

# Check CUDA availability
print(torch.cuda.is_available())
print(torch.cuda.device_count())
print(torch.cuda.get_device_name(0))

# Move tensors to GPU
device = torch.device('cuda')
x = torch.randn(1000, 1000).to(device)
y = torch.randn(1000, 1000).to(device)
z = x @ y  # Matrix multiplication on GPU

# Move model to GPU
model = MyModel().to(device)

# DataLoader for efficient loading
train_loader = DataLoader(
    train_dataset,
    batch_size=32,
    shuffle=True,
    num_workers=4,  # Parallel data loading
    pin_memory=True  # Faster GPU transfer
)

# Mixed precision training
scaler = torch.cuda.amp.GradScaler()
with torch.cuda.amp.autocast():
    output = model(data)
    loss = criterion(output, target)

scaler.scale(loss).backward()
scaler.step(optimizer)
scaler.update()
\end{minted}
\end{code}

\begin{properties}[CUDA Principles]
\begin{itemize}
    \item GPUs excel at parallel operations (matrix mult, convolutions)
    \item Data transfer between CPU and GPU is expensive
    \item Batching improves GPU utilization
    \item Mixed precision (FP16) reduces memory and increases speed
\end{itemize}
\end{properties}

\section{Chapter Summary}

\begin{enumerate}
    \item Training involves iterative optimization over batches
    \item Hyperparameters (LR, batch size) critically affect training
    \item Regularization prevents overfitting
    \item Backpropagation computes gradients efficiently
    \item PyTorch provides flexible deep learning primitives
    \item CUDA enables GPU acceleration for training
\end{enumerate}
